\subsection{逐点排除法}
\label{subsec:lab4-point-exclusion}

\paragraph{基本原理}

逐点排除法从所有八位二进制向量出发，对每个非法点建立一条只排除该点的不等式。将这些不等式同时加入模型后，所有非法点都被排除，而合法点全部保留。

\paragraph{数学建模}

设非法点集为
\[
U=\{0,1\}^{8}\setminus V.
\]
对于任意$u=(u_1,\ldots,u_8)\in U$，定义$z$与$u$的Hamming距离
\[
d_u(z)
=\sum_{j:u_j=0}z_j+\sum_{j:u_j=1}(1-z_j).
\]
当$u_j=0$时，$z_j$恰好表示该位是否不同；当$u_j=1$时，$1-z_j$恰好表示该位是否不同。因此，在二进制变量上，$d_u(z)$等于两向量不同坐标的数量。

要求$z$至少有一位不同于每个非法点，可得模型
\[
\boxed{
\begin{aligned}
&\sum_{j:u_j=0}z_j+\sum_{j:u_j=1}(1-z_j)\ge1,
&&u\in U,\\
&z_j\in\{0,1\},
&&j=1,\ldots,8.
\end{aligned}
}
\]
由于$u$是已知的常量向量，每条不等式也可以整理为标准线性形式
\[
\sum_{j:u_j=0}z_j-\sum_{j:u_j=1}z_j
\ge1-\sum_{j=1}^{8}u_j,
\qquad u\in U.
\]

例如，输入为$0$、输出为$1$对应的点为$u=0000\,0001$。对于本实验的置换S盒，该点既不属于差分合法集合，也不属于线性合法集合，其排除约束为
\[
z_1+z_2+z_3+z_4+z_5+z_6+z_7+(1-z_8)\ge1.
\]
将$0000\,0001$代入，左端为$0$；改变任意一位，左端便至少为$1$。这说明约束由非法点的坐标直接生成，不需要求解器先寻找其系数。

\paragraph{正确性分析}

对于固定的非法点$u$，若$z=u$，则$d_u(z)=0$，违反对应约束；若$z\ne u$，则至少一个坐标不同，从而$d_u(z)\ge1$。因此，一条约束在二进制空间中恰好排除其对应的非法点。

若$z\in V$，则$z$不同于每个$u\in U$，故满足全部不等式。若$z\notin V$，则$z\in U$，模型中存在针对$u=z$的约束，该约束必然被违反。由此可知，模型的二进制可行点集恰好为$V$。

\paragraph{用于差分与线性传播}

差分建模时，将DDT零位置对应的八位向量作为$U$；线性建模时，将LAT零位置对应的八位向量作为$U$。两种情况下均只使用原始变量$z$，不需要额外的选择变量。
